\chapter{Požadavky na výukovou aplikaci}
\label{chap:requirements}

Pro úspěšný vývoj a~implementaci webové aplikace pro vizualizaci datových struktur a~algoritmů je nezbytné stanovit jasné požadavky, které budou sloužit jako vodítko během celého vývojového procesu.
Jako inspirace pro tyto požadavky byly použity existující nástroje popsané v~\customref{Kapitole}{chap:other_apps} a~také obecné principy návrhu softwaru a~uživatelského rozhraní.
Požadavky jsou rozděleny do kategorií dle FURPS analýzy, zahrnujících funkcionalitu, použitelnost, spolehlivost, výkonnost a~podporovatelnost a~jsou popsány v~následujících kapitolách.

\section{Funkcionalita}
\label{sec:functionality}

Tato kapitola definuje hlavní funkční požadavky na systém, které zahrnují vizualizační schopnosti, interaktivní prvky a~vzdělávací obsah.

\begin{enumerate}
    \item \textbf{Moduly pro vizualizaci datových struktur}
          \begin{itemize}
              \item Implementace a~vizualizace struktur: zásobník, fronta, lineární seznam, halda a~vyhledávací stromy.
              \item Grafická simulace standardních operací nad těmito strukturami.
          \end{itemize}

    \item \textbf{Moduly pro simulaci algoritmů}
          \begin{itemize}
              \item Algoritmy řazení: Bubble sort, Insertion sort, Selection sort, Merge sort, Quick sort a~Heap sort.
              \item Algoritmy vyhledávání: lineární a~binární vyhledávání.
          \end{itemize}

    \item \textbf{Interaktivní ovládání simulací}
          \begin{itemize}
              \item Funkce krokování algoritmu (vpřed i~vzad) pro detailní analýzu průběhu.
              \item Možnost definice vstupních dat v~reálném čase s~okamžitým promítnutím do vizualizace.
          \end{itemize}

    \item \textbf{Teoretický a~metodický podklad}
          \begin{itemize}
              \item Formální vysvětlení principů a~typických scénářů užití u~každé struktury a~algoritmu.
              \item Specifikace časové a~prostorové algoritmické složitosti.
              \item Zobrazení odpovídajícího pseudokódu algoritmů pro lepší pochopení implementace.
          \end{itemize}

    \item \textbf{Lokalizace rozhraní}
          \begin{itemize}
              \item Plná podpora vícejazyčného rozhraní v~rozsahu českého a~anglického jazyka.
          \end{itemize}

    \item \textbf{Architektura na straně klienta}
          \begin{itemize}
              \item Provoz aplikace v~rámci webového prohlížeče bez nutnosti instalace dodatečných doplňků.
              \item Provádění veškerých simulací a~interakcí lokálně na zařízení uživatele.
              \item Persistence uživatelského nastavení prostřednictvím webového úložiště (local storage).
          \end{itemize}
\end{enumerate}

\section{Použitelnost}
\label{sec:usability}

Cílem požadavků na použitelnost je zajistit nízkou bariéru vstupu a~intuitivní ovládání pro cílovou skupinu uživatelů.

\begin{enumerate}
    \item \textbf{Ergonomie a~ovládání}
          \begin{itemize}
              \item Intuitivní uživatelské rozhraní doplněné o~kontextové nápovědy pro složitější ovládací prvky.
              \item Konzistentní navigační systém pro přehledný přechod mezi moduly struktur a~algoritmů.
          \end{itemize}

    \item \textbf{Didaktická podpora}
          \begin{itemize}
              \item Synergie vizualizace, teorie a~vysvětlujících textů pro podporu uživatelů s~různou úrovní vstupních znalostí.
          \end{itemize}

    \item \textbf{Vizuální prezentace}
          \begin{itemize}
              \item Implementace plynulých animací pro jasnou identifikaci změn stavu datových struktur nebo prvků v~poli.
              \item Moderní, přehledný a~čistý design uživatelského rozhraní odpovídající současným standardům (UI/UX).
          \end{itemize}
\end{enumerate}

\section{Spolehlivost}
\label{sec:reliability}

Požadavky na spolehlivost se zaměřují na stabilitu systému a~korektní zpracování uživatelských dat.

\begin{enumerate}
    \item \textbf{Stabilita systému}
          \begin{itemize}
              \item Zajištění stabilního a~konzistentního chování aplikace při zpracování datových sad o~rozsahu až 100 prvků.
              \item Udržení korektního stavu aplikace i~při opakovaných uživatelských interakcích.
          \end{itemize}

    \item \textbf{Prevence nevalidních operací}
          \begin{itemize}
              \item Aktivní blokování operací, které nejsou v~daném stavu aplikace přípustné.
              \item Vizuální zpětná vazba pro nedostupné ovládací prvky a~nevalidní vstupy, aby uživatel nemohl vyvolat chybovou akci.
          \end{itemize}
\end{enumerate}

\section{Výkonnost}
\label{sec:performance}

Důležitým aspektem výkonnosti je plynulost grafického výstupu a~rychlá odezva na interakci.

\begin{enumerate}
    \item \textbf{Optimalizace vykreslování}
          \begin{itemize}
              \item Zajištění stabilní plynulosti animací bez ohledu na objem vstupních dat (v~rámci stanovených limitů).
          \end{itemize}

    \item \textbf{Efektivita výpočtů}
          \begin{itemize}
              \item Optimalizované zpracování algoritmů na straně klienta pro minimalizaci latence mezi uživatelským vstupem a~vizuální odezvou.
          \end{itemize}
\end{enumerate}

\section{Podporovatelnost a~udržitelnost}
\label{sec:supportability}

Požadavky na podporovatelnost a~udržitelnost definují technický rámec pro budoucí rozšiřování a~údržbu aplikace.

\begin{enumerate}
    \item \textbf{Modularita a~separace zájmů}
          \begin{itemize}
              \item Striktní oddělení vizualizační vrstvy od aplikační logiky.
              \item Definice jasného a~jednoduchého rozhraní (API) mezi jednotlivými komponentami systému.
              \item Architektura UI umožňující snadné úpravy vzhledu bez zásahu do logiky (např. pomocí CSS proměnných).
          \end{itemize}

    \item \textbf{Použití moderních technologií}
          \begin{itemize}
              \item Vývoj klientské části využívající komponentní framework Svelte.
              \item Využití specializovaných knihoven pro vizualizaci dat (např. D3.js nebo vis.js).
          \end{itemize}

    \item \textbf{Dokumentace a~rozšiřitelnost}
          \begin{itemize}
              \item Detailní dokumentace systémové architektury pro usnadnění budoucího vývoje.
              \item Návrh systému umožňující snadné přidávání nových datových struktur, algoritmů nebo nových podporovaných jazyků.
          \end{itemize}

    \item \textbf{Správa kódu a~verzování}
          \begin{itemize}
              \item Správa zdrojových kódů a~historických změn prostřednictvím verzovacího systému Git.
          \end{itemize}
\end{enumerate}
